\documentclass[11pt]{article}
\usepackage[margin=.85in]{geometry}
\usepackage{amsmath,amssymb,amsthm,booktabs,array}
\usepackage[T1]{fontenc}
\usepackage{lmodern}
\usepackage[colorlinks=true,urlcolor=blue,citecolor=blue,linkcolor=blue,breaklinks=true]{hyperref}
\newtheorem{theorem}{Theorem}
\newtheorem{lemma}[theorem]{Lemma}
\newtheorem{corollary}[theorem]{Corollary}
\newcommand{\Z}{\mathbb Z}
\newcommand{\N}{\mathbb N}
\newcommand{\T}{\mathcal T}
\newcommand{\Hh}{\mathcal H}
\title{Proof of Barker's tea-tiling conjectures\\OEIS A299285--A299286, FORMULA\\(February 11, 2018)}
\author{Haobo Ma\\\small ChronoAI Pte Ltd, Singapore
\and Wenlin Zhang\\\small National University of Singapore\\\small ChronoAI Pte Ltd, Singapore}
\date{November 8, 2026}
\begin{document}
\maketitle
\begin{abstract}
We prove Colin Barker's generating-function and recurrence conjectures of
February 11, 2018 in the FORMULA fields of OEIS A299285 and A299286
for the tea uniform tiling.
An exact distance formula reduces its graph spheres to integer hexagons in
triangular layers. Their counts give a quadratic quasipolynomial of least
period four, establish Stefano Spezia's conditional formula of June 8, 2024,
and give the corresponding cubic ball count. We also prove that the proposed
recurrences have the least possible eventual orders, six and seven. The same
layer count proves the companion tcd formulas in A299287 and A299288.
For arbitrary orientations of the tetrahedron--octahedron slabs, the tcd
and tea counts are sharp lower and upper bounds, with sphere-count gap
$\lfloor n^2/8\rfloor$.
\end{abstract}

\section{The conjectures and their graph}
For a vertex in the edge graph of the tea uniform tiling, let $a(n)$ count
vertices at graph distance exactly $n$, and put $b(n)=\sum_{m=0}^n a(m)$.
Barker's formulas in the FORMULA fields of A299285 and A299286
\cite{shell,ball} predict the generating functions below. Sloane's
June 12, 2024 FORMULA annotation in A299285 says,
``I suspect Barker's formulas only conjectures.'' Spezia's June 8, 2024
FORMULA contribution begins ``If the above formulas are true''; its
closed form is consequently a conditional deduction, rather than a
graph-counting proof. The paired A299286 entry contains Barker's
generating function and recurrence with the same February 11, 2018 date.
The tiling is Gr\"unbaum's tiling~4: triangular-prism slabs alternate with
reflected tetrahedron--octahedron slabs \cite[p.~54, List~1, no.~4]{grunbaum}.

Geometric counting proofs for planar coordination sequences are given by
Goodman-Strauss and Sloane \cite{coloring}. Inoue and Nakamura develop
Ehrhart-theoretic methods for periodic graphs \cite{inoue1,inoue2}; their
second paper gives a general algorithm for obtaining explicit growth
series. Its three-dimensional examples are the carbon allotropes K6 and
CFS, rather than the two tilings considered here. For tea and tcd, the
proof below gives exact distances and counts directly, and also applies to
nonperiodic choices of slab orientations.

Here is an integer coordinate model, specifying the graph completely.
Its vertices are $(i,j,k)\in\Z^3$. Within a layer, the displacements are
$\pm(1,0,0)$, $\pm(0,1,0)$ and $\pm(1,-1,0)$.
Write $e_1=(1,0)$, $e_2=(0,1)$ and $\T=\{0,e_1,e_2\}$.
Between layers $k$ and $k+1$, the allowed upper-minus-lower planar
displacements are
\begin{equation}\label{edges}
\begin{array}{c|cccc}
k\bmod4&0&1&2&3\\\hline
\text{displacements}&\T&\{0\}&-\T&\{0\}.
\end{array}
\end{equation}
All edges are undirected. To identify the geometry explicitly, embed $e_1,e_2$
as unit vectors at $60^\circ$ and put $\delta=(e_1+e_2)/3$.
The successive layer offsets, modulo four, are $0,-\delta,-\delta,0$;
the gaps are $\sqrt{2/3},1,\sqrt{2/3},1$.
The three edges across a tetrahedron--octahedron slab have planar offsets
$-\delta,e_1-\delta,e_2-\delta$, or their negatives; across a prism slab
there is just the vertical edge. These are exactly \eqref{edges} for the
reflected slab construction in tiling~4.
Planar translations, together with
$(i,j,k)\mapsto(-i,-j,k+2)$ and
$(i,j,k)\mapsto(-i,-j,1-k)$, preserve \eqref{edges} and act transitively
on vertices. It therefore suffices to count about $(0,0,0)$.

\begin{theorem}\label{main}
Let $(c_0,c_1,c_2,c_3)=(16,17,12,17)$ and
$(d_0,d_1,d_2,d_3)=(32,38,24,30)$. Then $a(0)=1$ and, for $n\ge1$,
\begin{equation}\label{shellformula}
a(n)=\frac{63n^2+c_{n\bmod4}}8.
\end{equation}
For every $n\ge0$,
\begin{equation}\label{ballformula}
b(n)=\frac{84n^3+126n^2+104n+d_{n\bmod4}}{32}.
\end{equation}
With $P(x)=1+8x+14x^2+17x^3+14x^4+8x^5+x^6$,
\begin{equation}\label{gf}
\sum_{n\ge0}a(n)x^n=\frac{P(x)}{(1-x)^3(1+x)(1+x^2)},\qquad
\sum_{n\ge0}b(n)x^n=\frac{P(x)}{(1-x)^4(1+x)(1+x^2)}.
\end{equation}
\end{theorem}

\section{An exact distance formula}
For nonnegative integers $p,q$, define the integer hexagon
\[
\Hh(p,q)=\{(i,j)\in\Z^2:-q\le i,j,i+j\le p\}.
\]
The set of sums of $p$ members of $\T$ and $q$ members of $-\T$ is
exactly $\Hh(p,q)$. Indeed its elements satisfy the six inequalities.
Conversely, for a point in the hexagon the sum of its positive coordinates
is at most $p$, and the sum of the absolute values of its negative
coordinates is at most $q$. Taking the corresponding copies of
$e_1,e_2,-e_1,-e_2$, and padding with zeros, gives the required sum.

Let $p_k,q_k$ count the positive and negative special displacements in a
monotone vertical path from layer zero to layer $k$. From \eqref{edges},
\begin{equation}\label{slabs}
(p_k,q_k)=
\begin{cases}
(\lfloor(k+3)/4\rfloor,\lfloor(k+1)/4\rfloor),&k\ge0,\\
(\lfloor(-k+2)/4\rfloor,\lfloor(-k)/4\rfloor),&k<0.
\end{cases}
\end{equation}
Thus the planar endpoints attainable in $|k|$ vertical steps are precisely
$\Hh(p_k,q_k)$. Define
\begin{align*}
R_{p,q}(i,j)=\max\{&0,i-p,j-p,i+j-p,\\
                  &-q-i,-q-j,-q-i-j\}.
\end{align*}
\begin{lemma}\label{metric}
For every $(i,j,k)\in\Z^3$ the distance from the origin is
\[
D(i,j,k)=|k|+R_{p_k,q_k}(i,j).
\]
\end{lemma}
\begin{proof}
First, the triangular-lattice distance of $(i,j)$ from $\Hh(p,q)$ equals
$R=R_{p,q}(i,j)$. Each horizontal edge changes each of the six affine
expressions defining $R$ by at most one, so at least $R$ steps are needed.
If $R>0$, a step towards zero reduces it by exactly one: for $i,j\ge0$
decrease a positive coordinate; for $i,j\le0$ increase a negative
coordinate; for $i>0>j$ use $(-1,1)$, and for $j>0>i$ use $(1,-1)$.
In the first two cases $R$ is respectively
$\max(0,i+j-p)$ and $\max(0,-i-j-q)$.
In the third case it is $\max(0,i-p,-j-q)$, and the fourth is symmetric.
These formulas verify the asserted decrease, including axis points.
After $R$ steps the point lies in the hexagon.
Reverse that horizontal path after a monotone vertical path to its endpoint;
this constructs a path of length $D$.

For the reverse inequality, $D$ is one-Lipschitz on every edge.
This holds horizontally by the preceding affine-expression bound.
Orient a vertical edge away from layer zero. At a prism slab its planar
coordinates and $p,q$ are unchanged. At a positive special slab its new
data are $(z+t,p+1,q)$ with $t\in\T$.
Each expression in the defining maximum for the new residual is either
the corresponding old expression or that expression minus one;
the constant zero stays zero. Hence $R-1\le R'\le R$.
The negative special slab has data $(z-t,p,q+1)$ and the same bound.
Since $|k|$ increases by one, $0\le D'-D\le1$ in either case.
Finally $D(0,0,0)=0$, so every path from the origin has length at least
$D$. This proves the formula.
\end{proof}

\section{Counting layers}
\begin{lemma}\label{hexcount}
For $p,q\ge0$,
\begin{equation}\label{hex}
|\Hh(p,q)|=C(p,q):=\frac{p^2+4pq+q^2+3p+3q+2}{2}.
\end{equation}
The number of planar points with residual exactly $r$ is $C(p,q)$ for
$r=0$, and $3(p+q)+6r$ for $r\ge1$.
\end{lemma}
\begin{proof}
For each $-q\le i\le-1$ the allowed row has $p+q+1+i$ points;
for $0\le i\le p$ it has $p+q+1-i$ points.
Summing these arithmetic progressions gives \eqref{hex}.
For $r\ge0$, the residual is at most $r$ exactly on
$\Hh(p+r,q+r)$. For $r\ge1$, subtract \eqref{hex} at
$(p+r-1,q+r-1)$ from its value at $(p+r,q+r)$.
\end{proof}

\begin{proof}[Proof of Theorem~\ref{main}]
For $n\ge1$, Lemmas~\ref{metric} and \ref{hexcount} give
\begin{equation}\label{layersum}
a(n)=C(p_n,q_n)+C(p_{-n},q_{-n})+
\sum_{k=-n+1}^{n-1}\bigl(3(p_k+q_k)+6(n-|k|)\bigr).
\end{equation}
The floor formulas imply
$p_m+q_m+p_{-m}+q_{-m}=m$ for $m\ge1$.
Also $p_0=q_0=0$ and
$\sum_{k=-n+1}^{n-1}(n-|k|)=n^2$.
Writing
\[
Q_n=p_n^2+4p_nq_n+q_n^2+p_{-n}^2+4p_{-n}q_{-n}+q_{-n}^2,
\]
equation~\eqref{layersum} becomes
\[
a(n)=\frac{15n^2+Q_n+4}{2}.
\]
For $n=4t+s$, the endpoint parameters and $Q_n$ are
\[
\begin{array}{c|c|c|c}
s&(p_n,q_n)&(p_{-n},q_{-n})&Q_n\\\hline
0&(t,t)&(t,t)&12t^2\\
1&(t+1,t)&(t,t)&12t^2+6t+1\\
2&(t+1,t)&(t+1,t)&12t^2+12t+2\\
3&(t+1,t+1)&(t+1,t)&12t^2+18t+7
\end{array}
\]
Substitution proves \eqref{shellformula}; the radius-zero sphere has one
vertex. If $u(n)=\cos(\pi n/2)\in\{1,0,-1,0\}$ by residue class, the
same formula is
\begin{equation}\label{spezia}
a(n)=\frac{31-3(-1)^n+126n^2+4u(n)}{16}\quad(n\ge1).
\end{equation}
This is Spezia's conditional formula in A299285, with $u(n)=$ A056594$(n)$~\cite{periodic}.
Summing \eqref{spezia}, using
$\sum_{m=1}^n u(m)=0,0,-1,-1$ for $n\equiv0,1,2,3\pmod4$,
gives \eqref{ballformula}.
Finally,
\[
\sum_{n\ge0}a(n)x^n=1+\frac1{16}\left(
\frac{126x(1+x)}{(1-x)^3}+\frac{31x}{1-x}
+\frac{3x}{1+x}+\frac4{1+x^2}-4\right).
\]
Combining fractions gives the first identity in \eqref{gf}.
Partial sums divide the generating function by $1-x$, giving the second.
\end{proof}

\begin{corollary}\label{recurrences}
The following recurrences hold for $n\ge7$:
\begin{align*}
a(n)&=2a(n-1)-a(n-2)+a(n-4)-2a(n-5)+a(n-6),\\
b(n)&=3b(n-1)-3b(n-2)+b(n-3)+b(n-4)\\
&\hspace{12mm}-3b(n-5)+3b(n-6)-b(n-7).
\end{align*}
Their least eventual orders over $\mathbb Q$ are six and seven,
respectively. The least quasipolynomial period is four for each sequence
(for $a(n)$ on $n\ge1$).
\end{corollary}
\begin{proof}
Multiply \eqref{gf} by its denominators and compare coefficients beyond
degree six. The numerators do not cancel any denominator factor:
$P(1)=63$, $P(-1)=-3$, and $P(x)\equiv-x\pmod{x^2+1}$.
An eventual recurrence with polynomial $E(x)$ makes the generating
function times $E(x)$ a polynomial; therefore the reduced denominator
divides $E(x)$. Its degree is the minimum possible order.
Similarly, an eventual quasipolynomial of period $h$ has a generating
function whose poles are $h$th roots of unity, by summing each polynomial
residue class separately. The uncancelled poles at $i,-i$ force $4\mid h$;
the displayed formulas supply period four.
\end{proof}

\section{The cross-referenced tcd tiling}
\begin{corollary}\label{tcdresults}
The formulas in the FORMULA fields of A299287 and A299288
dated February 11, 2018, and Spezia's June 8, 2024 exponential
generating function in A299287, hold for the tcd graph. Its sphere and
ball recurrences have least eventual orders four and five; both least
quasipolynomial periods are two.
\end{corollary}
\begin{proof}
The same argument gives the formulas in A299287 and A299288
\cite{tcd,tcdball}, concerning Gr\"unbaum's tiling~3. Its tetrahedron--octahedron
slabs have the same orientation rather than alternating orientations:
replace both even columns of \eqref{edges} by $\T$.
Consequently, for $k\ge0$ its parameters are
$(p_k,q_k)=(\lceil k/2\rceil,0)$ and
$(p_{-k},q_{-k})=(0,\lfloor k/2\rfloor)$.
The distance proof and \eqref{layersum} apply without change; in particular
the same $3(p+q)+6r$ counts the nonzero residuals.
Writing $\alpha(n)$ for its spheres and $\beta(n)$ for its balls gives
$Q_n=\lceil n/2\rceil^2+\lfloor n/2\rfloor^2$ and hence
\[
\alpha(0)=1,\qquad
\alpha(n)=\frac{31n^2+8+\epsilon_n}{4}\ (n\ge1),\qquad
\beta(n)=\frac{62n^3+93n^2+82n+24+3\epsilon_n}{24}\ (n\ge0),
\]
where $\epsilon_n$ is zero for even $n$ and one for odd $n$.
Summing the quadratic formula, or its generating series, gives
\[
\sum_{n\ge0}\alpha(n)x^n=\frac{1+8x+13x^2+8x^3+x^4}{(1-x)^3(1+x)},\qquad
\sum_{n\ge0}\beta(n)x^n=\frac{1+8x+13x^2+8x^3+x^4}{(1-x)^4(1+x)}.
\]
These are Barker's formulas and imply, for $n\ge5$,
\begin{align*}
\alpha(n)&=2\alpha(n-1)-2\alpha(n-3)+\alpha(n-4),\\
\beta(n)&=3\beta(n-1)-2\beta(n-2)-2\beta(n-3)+3\beta(n-4)-\beta(n-5).
\end{align*}
The numerator takes values $31$ and $-1$ at $1$ and $-1$, so these
recurrences have least eventual orders four and five, and each
quasipolynomial has least period two, by the pole argument above.
Finally Spezia's June 8, 2024 exponential generating function in A299287 is
\[
\sum_{n\ge0}\alpha(n)\frac{x^n}{n!}
=\frac{(8+31x+31x^2)\cosh x+(9+31x+31x^2)\sinh x-4}{4}.
\]
Indeed apply $31(x\,d/dx)^2$ to $\cosh x+\sinh x$, add
$8\cosh x+9\sinh x$, divide by four, and subtract one to correct
the constant term. This proves all formulas displayed in these two
companion entries.
\end{proof}

\section{Arbitrary slab orientations}
Keep the triangular layers and the prism slabs at odd $k$ in
\eqref{edges}, but independently choose $\T$ or $-\T$ at each even
boundary $k$. This defines a graph $G_\sigma$ for every two-sided sign
sequence $\sigma$, without a periodicity hypothesis. Let $a_\sigma(n)$
count its sphere about $(0,0,0)$ and $b_\sigma(n)$ its ball. For a
monotone path to layer $k$, let $p_k,q_k$ count the $\T,-\T$
displacements respectively, reversing signs on downward crossings.

\begin{corollary}\label{arbitrary}
For $n\ge1$, put $r=\lceil n/2\rceil$ and $s=\lfloor n/2\rfloor$.
Then
\[
a_\sigma(n)=\alpha(n)+p_nq_n+p_{-n}q_{-n},\qquad
\alpha(n)\le a_\sigma(n)\le a(n).
\]
Both bounds hold for every $n\ge0$, simultaneously, with the lower
bound attained by tcd and the upper bound by tea. The balls satisfy
$\beta(n)\le b_\sigma(n)\le b(n)$. In particular,
\[
a(n)-\alpha(n)=\left\lfloor\frac{n^2}{8}\right\rfloor,
\qquad
b(n)-\beta(n)=\sum_{m=1}^n\left\lfloor\frac{m^2}{8}\right\rfloor.
\]
For $n\ge1$, equality in the lower sphere bound holds exactly when
$p_nq_n=p_{-n}q_{-n}=0$. Equality in the upper bound holds exactly when
$|p_n-q_n|\le1$ and $|p_{-n}-q_{-n}|\le1$.
\end{corollary}
\begin{proof}
The distance proof uses only that an outward special crossing adds
one positive or negative slab. It therefore holds for $G_\sigma$.
The endpoint identities are $p_n+q_n=r$, $p_{-n}+q_{-n}=s$.
Also $p_m+q_m+p_{-m}+q_{-m}=m$ for every $m\ge1$.
The layer sum consequently remains
\[
a_\sigma(n)=\frac{15n^2+Q_n+4}{2},\quad
Q_n=r^2+s^2+2(p_nq_n+p_{-n}q_{-n}).
\]
The first term is $\alpha(n)=(15n^2+r^2+s^2+4)/2$.
For nonnegative integers of sum $r$, the product ranges from zero to
$\lfloor r^2/4\rfloor$, with maximum exactly when their difference is
at most one; the same holds for $s$. Now
\[
\left\lfloor\frac{r^2}{4}\right\rfloor+
\left\lfloor\frac{s^2}{4}\right\rfloor
=\left\lfloor\frac{n^2}{8}\right\rfloor.
\]
For $n=4t,4t+1,4t+2,4t+3$, both sides are respectively
$2t^2,2t^2+t,2t^2+2t,2t^2+3t+1$.
The tcd parameters put each product at zero; the tea parameters in
\eqref{slabs} balance each pair and attain both maxima. This proves
the inequalities, the equality conditions and the exact gap.
All radius-zero counts are one. Summing proves the ball statements.
\end{proof}
This extension is not posed in the cited entries. It concerns the
specified layered graphs and does not assert vertex transitivity for
an arbitrary sign sequence.

\section{Coverage and verification}
\noindent\begin{tabular}{@{}p{.48\linewidth}p{.47\linewidth}@{}}
\toprule
Source item & Mathematical coverage\\\midrule
A299285, Barker, Feb.~11, 2018: generating function and recurrence
& Proved: Theorem~\ref{main}, Corollary~\ref{recurrences}.\\[3pt]
A299285, Spezia, June~8, 2024: conditional closed form
& Proved for the graph: equation~\eqref{spezia}; the prior conditional algebraic implication is credited to Spezia.\\[3pt]
A299286, Barker, Feb.~11, 2018: partial-sum generating function and recurrence
& Proved: Theorem~\ref{main}, Corollary~\ref{recurrences}.\\[3pt]
A056594: period-four cosine sequence
& Its defining values identify $u(n)$ in \eqref{spezia}.\\[3pt]
A299287, FORMULA, Barker, Feb.~11, 2018: closed forms, generating function, recurrence
& Proved: Corollary~\ref{tcdresults}, including least order four and least period two.\\[3pt]
A299287, FORMULA, Spezia, June~8, 2024: exponential generating function
& Proved: Corollary~\ref{tcdresults}, by summing the sphere formula.\\[3pt]
A299288, FORMULA, Barker, Feb.~11, 2018: ball formulas, generating function, recurrence
& Proved: Corollary~\ref{tcdresults}, including least order five and least period two.\\[3pt]
A299285--A299288, DATA and finite sequence programs
& Finite checks agree; the all-radius statements have the written proofs above.\\[3pt]
Arbitrary signs, sharp bounds, equality cases and tea--tcd gap
& Not posed; proved in Corollary~\ref{arbitrary}.\\[3pt]
Other tilings in the common OEIS cross-reference list
& Not covered by this note.\\\bottomrule
\end{tabular}

\medskip
The pinned Lean file
\href{https://github.com/the-omega-institute/trureturing/blob/e7f7c25c7de35e89418be730f50ec8fb94949133/D5/S3/StatisticalMechanics/TeaCoordinationSequence.lean}{\texttt{TeaCoordinationSequence.lean}}
defines \texttt{teaGraph} by \eqref{edges} and \texttt{shell} as its
distance-sphere cardinality. Its declaration
\begin{quote}\small
\texttt{D5.S3.StatisticalMechanics.TeaCoordinationSequence.result}
\end{quote}
proves \texttt{claim}: the four identities
$8a(n)=63n^2+c_{n\bmod4}$ for $n>0$.
The identification with Gr\"unbaum's tiling, the generating functions,
the partial sums, the least-order and least-period assertions, and all tcd
results, and the arbitrary-orientation extension have the written proofs above.
The accompanying \texttt{verify.py} supplies
finite integer checks for both graphs: breadth-first spheres through radius 32,
the distance formula at every vertex in each ball, the tabulated OEIS
values, and layer sums and generating-function coefficients through 1000.
For arbitrary orientations, it checks all 256 assignments on the eight even
boundaries $-8,-6,\ldots,6$, with graph distances and bounds through radius
eight. It also checks the tea--tcd gaps through radius 1000.
Its JSON output states the finite parameter ranges.

\begin{thebibliography}{9}
\bibitem{shell} N.~J.~A. Sloane, OEIS A299285,
\emph{Coordination sequence for ``tea'' 3D uniform tiling}, created
February 10, 2018; formulas by Colin Barker, February 11, 2018, and
Stefano Spezia, June 8, 2024; Sloane annotation June 12, 2024.
\url{https://oeis.org/A299285}. Revision 25, consulted November 8, 2026.
\bibitem{ball} N.~J.~A. Sloane, OEIS A299286,
\emph{Partial sums of A299285}, created February 10, 2018;
conjectures by Colin Barker, February 11, 2018.
\url{https://oeis.org/A299286}. Revision 9, consulted November 8, 2026.
\bibitem{periodic} W. Lang, OEIS A056594,
\emph{Period 4: repeat [1,0,-1,0]; expansion of 1/(1+x\textasciicircum 2)},
\url{https://oeis.org/A056594}. Revision 187, consulted November 8, 2026.
\bibitem{tcd} N.~J.~A. Sloane, OEIS A299287,
\emph{Coordination sequence for ``tcd'' 3D uniform tiling},
Barker formulas dated February 11, 2018; Spezia exponential generating function
dated June 8, 2024. \url{https://oeis.org/A299287}. Consulted November 8, 2026.
\bibitem{tcdball} N.~J.~A. Sloane, OEIS A299288,
\emph{Partial sums of A299287}, Barker formulas dated February 11, 2018.
\url{https://oeis.org/A299288}. Revision 12, consulted November 8, 2026.
\bibitem{grunbaum} B. Gr\"unbaum, \emph{Uniform tilings of 3-space},
Geombinatorics \textbf{4} (1994), 49--56.
\url{https://faculty.washington.edu/moishe/branko/BG199.Uniform\%20Tilings\%20of\%203-Space.pdf}.
\bibitem{coloring} C. Goodman-Strauss and N.~J.~A. Sloane,
\emph{A coloring-book approach to finding coordination sequences},
Acta Crystallographica A \textbf{75} (2019), 121--134.
\url{https://doi.org/10.1107/S2053273318014481}.
\url{https://arxiv.org/abs/1803.08530v3} (January 31, 2019).
\bibitem{inoue1} T. Inoue and Y. Nakamura,
\emph{Ehrhart theory on periodic graphs},
Algebraic Combinatorics \textbf{7} (2024), 969--1010.
\url{https://doi.org/10.5802/alco.367}.
\url{https://arxiv.org/abs/2305.08177v2} (April 9, 2024),
especially Section~5.3.
\bibitem{inoue2} T. Inoue and Y. Nakamura,
\emph{Ehrhart theory on periodic graphs II: Stratified Ehrhart ring theory},
Algebraic Combinatorics \textbf{8} (2025), 1193--1232.
\url{https://doi.org/10.5802/alco.445}.
\url{https://arxiv.org/abs/2310.19569v2} (July 3, 2024),
especially Sections~5.3--5.4.
\end{thebibliography}
\noindent\small\copyright\ 2026 Haobo Ma and Wenlin Zhang.
Licensed under Creative Commons Attribution 4.0 International (CC BY 4.0).
\url{https://creativecommons.org/licenses/by/4.0/}
\end{document}
